\item \label{itm:rad} {\bf [6 pts.] Thermal stability of cratons} \\
Cratons are typically viewed as regions that are thermally cool, which
contributes to their longevity/stability.  In this problem, you will explore the
effect of crustal radioactivity on crustal temperatures both for the present-day
Earth and at a time \SI{3}{Ga} when many of these stable regions were being first
created.

This problem is accompanied by a spreadsheet (see template online) that you will
need to complete as part of the exercises below.
\begin{enumerate}
    \item Compute the average heat production for a granitic and gabbroic
        compositions at the present day for crystallization ages of $<$\SI{200}{Ma},
        from 2000 to \SI{2800}{Ma}, and prior to \SI{2800}{Ma}.  The formulae should be
        directly entered into the worksheet labeled "heat production" and copied
        and modified as necessary for the other two rock ages.
        
        Note: the concentration of potassium is provided as K$_2$O, so you will
        need to convert to K.  And both K and U have multiple isotopes, so you
        will need to account for that in your calculation using the constants
        provided in the worksheet.

    \item In the same sheet, compute the average heat production for these same
        rock types and ages adjusted to a time 3 billion years ago (\SI{3}{Ga}).

    \item The values you previously computed will be automatically copied into
        the worksheet labeled ``temperature''.  We want to compute a set of
        geotherms in this new sheet for a simplified crust created for six
        crustal heat production scenarios corresponding to the different age
        rocks at both present day and at \SI{3}{Ga}.

        We will make some simplifying assumptions for a typically moderate
        crustal heat flow:
        \begin{itemize}
            \item upper crust is \SI{16}{km} thick; 
            \item lower crust is \SI{24}{km} thick;
            \item basal heat flow is \SI{36}{mW.m^{-2}}---use the conservation of
                energy to compute the heat flow at the top of each layer, all
                the way back to the surface; and
            \item thermal conductivity is temperature variable---the equations
                for thermal conductivity are already computed for you and will
                automatically update once you have computed the temperature.
        \end{itemize}

        If you enter the equations correctly into the first empty heat flow and
        temperature cells, you should be able to copy them to the end of each
        column.  With the appropriate \$ on cell columns/rows, you will also be
        able to copy the heat flow and temperature columns from the first
        scenario to all others.

    \item The remainder of the problem leads up to computing heat production for
        three hypothetical regions at a point \SI{3}{Ga} in the past so that the
        geotherms in the past.

        Start by determining the relationship between half-life and radioactive
        decay constants.  The concentration of a radioelement, P(t), is given by
        \begin{equation}
            P(t) = P_0 \exp (-\lambda t),
        \end{equation}
        where $P_0$ is the initial concentration, and $\lambda$ is the decay
        constant.  Hint: a half-life is defined as the time at which the ratio
        of the radioelement concentration is one-half the initial concentration.
        Show your work.

    \item Compute the heat production of for the heat production concentrations
        at a point in time \SI{3}{Ga} in the past.  To compute the heat production in
        the past, you will need to adjust the concentrations of the
        radioisotopes.

    \item Compute one-dimensional, steady-state geotherms for each of the
        situations above to the Moho at a time \SI{3}{Ga} in the past.

    \item Explain, in light of your computations, how changes in heat production
        concentrations of igneous rocks helps lead to stability of Archean
        lithosphere.
\end{enumerate}
\FloatBarrier
